\section*{الفصل \ref{ch:b3:virology} --- علم الفيروسات}

\begin{solution}{exo:b3:virology:1}
شلل الأطفال: IV، RNA موجب الطاق --- بوليميراز RNA معتمد على RNA.
والإنفلونزا: V، RNA سالب الطاق --- بوليميراز RNA خاص به، محمول في
\omterm{def:b3:virology:virus}{الفيريون}. وHIV: VI --- ناسخة عكسية (وإنتغراز). والهربس
البسيط: I، DNA مزدوج الطاق --- كان يمكن للمضيف مبدئيًا أن يوفر
كل شيء، لكنه يجلب بوليميراز DNA وكيناز ثيميدين خاصين به.
والتهاب الكبد B: VII --- ناسخة عكسية لنسخ RNA ما قبل الجينومي
إلى DNA. \omterm{def:b3:virology:virus}{والفيروس} العجلي: III، RNA مزدوج الطاق --- بوليميراز RNA معتمد على RNA
داخل الجسيم.
\end{solution}

\begin{solution}{exo:b3:virology:2}
دخل DNA العاثية ($^{32}$P) البكتيريا وبقي البروتين
($^{35}$S) في الخارج وأمكن قصّه؛ وورثت العاثيات الوليدة
$^{32}$P. وDNA إذن هو ما تحقنه العاثية وما
يوجّه صنع عاثيات جديدة. ولو كان البروتين يحمل التعليمات
لوُجد $^{35}$S في الداخل وانتقل إلى الذرية.
\end{solution}

\begin{solution}{exo:b3:virology:3}
$84/(0.1\times 10^{-7}) = 8.4\times 10^{9}$ وحدة تكوين لويحات في المليلتر.
فالمجهر يعدّ كل جسيم، بما فيه القفيصات الفارغة، والجسيمات
ذات الجينومات المعطوبة، والتي تخفق في الالتصاق أو الحقن؛ ولا تصنع
لويحات إلا الجسيمات التي تتم خمجًا.
\end{solution}

\begin{solution}{exo:b3:virology:4}
الالتصاق: مثبطات الدخول (المارافيروك)، والأجسام المضادة المعادِلة.
والدخول والاندماج: مثبطات الاندماج، وحاجبات تحميض الجسيمات الداخلية.
ونزع \omterm{def:b3:virology:virus}{الغلاف}: أدوية رابطة \omterm{def:b3:virology:virus}{للقفيصة} (البليكوناريل؛ واللينكابافير). والتعبير
والتضاعف: نظائر النوكليوزيد، ومثبطات الناسخة العكسية و
الإنتغراز، والإنترفيرون، \omterm{def:b3:rna-regulation:rnai}{والتداخل الرنائي}. والتجميع: مثبطات
البروتياز (فمتعددات الببتيد غير المشطورة لا تستطيع التجمع). والإطلاق:
مثبطات النورامينيداز (في الإنفلونزا)، والخلايا التائية السامة تقتل
الخلية قبل الإطلاق.
\end{solution}

\begin{solution}{exo:b3:virology:5}
البروتين الواحد من \qty{30}{kDa} هو $30\,000/110 = 273$ بقية، أي $818$
نوكليوتيد: أي \qty{0.82}{kb}، و\qty{18}{\%} من جينوم قدره \qty{4.5}{kb}.
وتشفير النسخ $180$ منفردةً كان ليقتضي \qty{147}{kb}،
أي ثلاثة وثلاثين ضعف الجينوم كله؛ وبروتين \omterm{def:b3:virology:virus}{القفيصة} نفسه
(\qty{5.4}{MDa}) يفوق الجينوم (\qty{1.5}{MDa}) أربعة أضعاف.
\end{solution}

\begin{solution}{exo:b3:virology:6}
$V = 10^{5}(3e^{-0.5t} - 0.5e^{-3t})/2.5$: فعند $t = 0.5$: $8.9\times
10^{4}$؛ وعند $t = 2$: $4.4\times 10^{4}$؛ وعند $t = 10$: $8\times 10^{2}$. ويهبط دون
$50$ حين $1.2\times 10^{5}e^{-0.5t} = 50$: أي $t = 2\ln(2400) \approx
\qty{16}{d}$.
\end{solution}

\begin{solution}{exo:b3:virology:7}
$uL = 0.36$؛ والنسبة الخالية من الخطأ $e^{-0.36} = 0.70$. وبثلاثة أضعاف: $uL =
1.08$، و$e^{-1.08} = 0.34$ --- أي إن معظم الذرية تحمل الآن طفرات و
يُمدَّد التسلسل الأصلح نحو العتبة. \omterm{def:b3:virology:virus}{وفيروس} كورونا عند
\qty{30}{kb} وعند $u = 10^{-4}$ كان ليكون عنده $uL = 3$ و\qty{5}{\%} فقط
من النسخ خالية من الخطأ، أي وراء العتبة؛ ونوكليازه الخارجي يجعل $u$
$10^{-5}$، فيصير $uL = 0.3$، أي ثلاثة أرباعها خالية من الخطأ.
\end{solution}

\begin{solution}{exo:b3:virology:8}
الانجراف: تتراكم الطفرات النقطية في الراصة الدموية والنورامينيداز
تحت انتقاء الأجسام المضادة، فتكون سلالة كل موسم
جديدة جزئيًا ومناعةُ العام الماضي عتيقة جزئيًا. والإزاحة: فالخمج المشترك
لخلية واحدة بسلالة بشرية وأخرى حيوانية يعيد توزيع القطع
الثماني، فتظهر زمرة راصة دموية من الطيور أو الخنازير في
\omterm{def:b3:virology:virus}{فيروس} متكيف مع البشر. ولا أحد عنده أجسام مضادة للزمرة الجديدة، فتكون
الجمهرة كلها قابلة دفعةً واحدة --- وهو شرط
الجائحة، وقد تحقق في 1918 و1957 و1968 و2009 براصة دموية جديدة في كل
مرة.
\end{solution}

\begin{solution}{exo:b3:virology:9}
كيناز الثيميدين في \omterm{def:b3:virology:virus}{فيروس} الهربس يفسفر الأسيكلوفير، وهو ما
لا تكاد تمسه كينازات الخلية؛ ثم تتم كينازات الخلية
ثلاثي الفوسفات، الذي يُدخله بوليميراز DNA الفيروسي، فينهي
السلسلة لخلوه من الهيدروكسيل 3$'$. والخلايا غير المخموجة لا تصنع الصورة
الفعالة قط، والبوليميراز الفيروسي يفضّلها على الخلوي.
والمقاومة: فقد كيناز الثيميدين الفيروسي أو تبدله (وهو
الأشيع)، أو طفرات في البوليميراز الفيروسي تستبعد
النظير.
\end{solution}

\begin{solution}{exo:b3:virology:10}
بإبقاء $T$ ثابتة تكون المعادلتان خطيتين في $(I, V)$ بالمصفوفة
$\begin{pmatrix} -\delta & \beta T \\ p & -c \end{pmatrix}$، وأثرها
$-(\delta + c) < 0$ ومحددها $\delta c - \beta Tp$. و
القيم الذاتية أثرها سالب، فتكون إحداها موجبة بالضبط حين يكون
المحدد سالبًا: أي $\beta Tp > c\delta$، وهذا يعني $R_{0} = \beta
Tp/(c\delta) > 1$. والعوامل: $p/\delta$ عدد الفيريونات التي تصنعها
خلية مخموجة في حياتها؛ و$\beta T/c$ عدد الخلايا التي يخمجها
\omterm{def:b3:virology:virus}{فيريون} في حياته؛ وجداؤهما عدد الخلايا المخموجة
التي تنتجها خلية مخموجة واحدة.
\end{solution}

\begin{solution}{exo:b3:virology:11}
$10^{6} = 2^{20}$: أي عشرون عمر نصف، و$880$ شهرًا، أي نحو $73$ سنة ---
أي أطول من عمر. والأدوية التي تحجب التضاعف لا تمس
فيروسًا مندمجًا لا يُستنسخ؛ فيدوم المستودع، وحين
يتوقف العلاج تعيد خلية واحدة مُعاد تنشيطها بدء الخمج. والشفاء
يقتضي إزالة المستودع أو إسكاته دائمًا: أي قتل الخلايا الكامنة
(بإعادة تنشيطها تحت العلاج فتموت أو تُقتل)، أو استئصال
\omterm{def:b3:virology:virus}{الفيروس} المندمج أو تعطيله، أو استبدال الجهاز المناعي بجهاز لا يستطيع \omterm{def:b3:virology:virus}{الفيروس}
إخماجه.
\end{solution}

\begin{solution}{exo:b3:virology:12}
$uL = 0.03$ طفرة لكل جينوم؛ و$10^{9}$ جينوم في اليوم تحمل $3\times
10^{7}$ طفرة بين $9\times 10^{4}$ طفرة مفردة ممكنة: فتنشأ كل
منها نحو $300$ مرة في اليوم. والطافر المزدوج بعينه ينشأ بمعدل
$10^{-12}$ لكل جينوم، أي $10^{-3}$ في اليوم --- أي مرة في ثلاث سنوات؛ و
الثلاثي عند $10^{-18}$، أي أبدًا. ودواءان بطفرتي مقاومة
مستقلتين يكفيان مبدئيًا، والثلاثة تعطي هامشًا لسوء
الالتزام. أما الدواءان اللذان يرتبطان بالإنزيم نفسه فقد تهزمهما طفرة
واحدة تبدّل الموقع أو انتظام الإنزيم، ولا تكون
طفرتا مقاومتهما مستقلتين؛ فيُعدّان أقل من اثنين.
\end{solution}

\begin{solution}{pb:b3:virology:1}
\textbf{1.} $2\times 10^{5}\times 1.5\times 10^{4} = 3\times 10^{9}$
\omterm{def:b3:virology:virus}{فيريون}.
\textbf{2.} المصفّى $cV = 3\times 3\times 10^{9} \approx 10^{10}$ في اليوم؛
والمنتَج مثله.
\textbf{3.} $I^{*} = cV^{*}/p = 9\times 10^{9}/500 = 1.8\times 10^{7}$
خلية مخموجة؛ ويموت منها $\delta I^{*} = 9\times 10^{6}$ في اليوم.
\textbf{4.} RNA: $2\times 9700\times 330 = 6.4\times 10^{6}$ دالتون؛ والبروتين
$2000\times 25\,000 = 5\times 10^{7}$ دالتون؛ والمجموع $5.6\times 10^{7}$ دالتون
$\approx 10^{-16}$ غرام. والفيريونات كلها معًا: $3\times 10^{9}\times
10^{-16} = 3\times 10^{-7}$ غرام، أي ثلث ميكروغرام --- أي أقل بثلاثة
آلاف مرة من حبة رمل.
\textbf{5.} $1.8\times 10^{7}/10^{12} = 2\times 10^{-5}$ من المخزون.
وعشر سنوات من $9\times 10^{6}$ حالة موت في اليوم هي $3\times 10^{10}$ خلية،
أي ثلاثة في المئة من المخزون: فالمخزون لا يخفق بطرح
حسابي بل لأن عقدًا من التجدد القسري والتنشيط
المزمن يستنفد القدرة على استبدالها.
\textbf{6.} $9\times 10^{6}/24 \approx 4\times 10^{5}$ جثة تُصفّى
في أي لحظة.
\textbf{7.} $V = 2\times 10^{5}(3e^{-0.5t} - 0.5e^{-3t})/2.5$: فعند $t = 1$:
$1.4\times 10^{5}$؛ وعند $t = 3$: $5.4\times 10^{4}$؛ وعند $t = 7$: $7\times
10^{3}$؛ وعند $t = 14$: $2\times 10^{2}$.
\textbf{8.} $2.4\times 10^{5}e^{-0.5t} = 50$: أي $t = 2\ln(4800) \approx 17$
يومًا.
\textbf{9.} الأيام 3--10 تقع على الطور البطيء: فميل $\ln V$ هناك هو
$-\delta$ (من $5.4\times 10^{4}$ إلى $1.6\times 10^{3}$ في سبعة أيام،
ومنه $\delta = 0.50$). أما الطور السريع فلا يدوم إلا ساعات ($1/c$)، فتحتاج $c$ إلى
عينات كل بضع ساعات في اليوم الأول، تُطابَق على صيغة الأسّيين.
\textbf{10.} الخلايا المخموجة كمونيًا، وهي تطلق \omterm{def:b3:virology:virus}{الفيروس} وهي
تُعاد تنشيطها؛ ومعدل تلاشيها $\ln 2/44$ في الشهر، أي نحو
$5\times 10^{-4}$ في اليوم.
\textbf{11.} \omterm{def:b3:virology:virus}{الفيريون} $\ln 2/3 = \qty{0.23}{d}$، أي نحو \qty{5.5}{h}؛ والخلية
المخموجة $\ln 2/0.5 = \qty{1.4}{d}$.
\textbf{12.} الحمل الثابت بدا فيروسًا لا يفعل شيئًا؛ أما
المعادلات فتُظهره حالة مستقرة يُصنع فيها $10^{10}$ \omterm{def:b3:virology:virus}{فيريون}
ويُصفّى وتُخمج $10^{7}$ خلية تائية وتُقتل كل يوم.
\textbf{13.} $9700\times 10^{-5} \approx 0.1$ طفرة لكل جينوم؛
و$10^{10}\times 0.1 = 10^{9}$ طفرة في اليوم.
\textbf{14.} $3\times 9700 \approx \num{29000}$ طفرة مفردة ممكنة؛
وتنشأ كل منها $10^{9}/\num{29000} \approx 3\times 10^{4}$ مرة في اليوم.
\textbf{15.} المزدوج بعينه: $10^{-10}$ لكل جينوم، أي نحو واحد في اليوم
بلا علاج؛ والثلاثي بعينه: $10^{-15}$، أي $10^{-5}$ في اليوم --- أي مرة في
ثلاثة قرون.
\textbf{16.} $10^{7}\times 10^{-15} = 10^{-8}$ في اليوم، أي $4\times 10^{-6}$ في
السنة.
\textbf{17.} بزوال دواء واحد، تكفي الطافرات المقاومة للآخرين:
فالمزدوج بعينه عند $10^{-10}$؛ و$7\times 10^{7}$ جينوم في
الأسبوع تعطي $7\times 10^{-3}$ متوقعًا --- وهو ما يزال غير مرجح، ما لم
يرتد الحمل نحو $10^{10}$ في اليوم أثناء الانقطاع، فيصير حينئذ
واحدًا في اليوم ويصير الإفلات مرجحًا. فالجرعات الفائتة هي كيف تصل المقاومة.
\textbf{18.} النظائر تشترك في الموقع الفعال للناسخة العكسية،
وطفرة واحدة هناك (أو طقم من طفرات نظائر الثيميدين) تبدّل
تعامل الإنزيم مع عدة منها: أي المقاومة المتصالبة. فالنظيران
النوكليوزيديان يُعدّان أقل من دواءين مستقلين، ولهذا تجمع النظم
بين الأصناف --- \omterm{def:b3:virology:drugs}{نظير نوكليوزيد}، ومثبط إنتغراز، \omterm{def:b3:virology:drugs}{ومثبط بروتياز}
أو غير نوكليوزيدي.
\textbf{19.} $20$ عمر نصف، و$880$ شهرًا، أي $73$ سنة؛ وإلى $10^{4}$،
$\log_{2}100 = 6.6$ عمر نصف، أي $290$ شهرًا، أو $24$ سنة.
\textbf{20.} من \omterm{def:b3:virology:virus}{فيريون} واحد إلى $3\times 10^{9}$ عند $R_{0} = 8$ لكل
جيل: $\ln(3\times 10^{9})/\ln 8 = 10.5$ جيلًا، أي ثلاثة أسابيع.
\textbf{21.} الصدم والقتل: أعد تنشيط الخلايا الكامنة تحت العلاج حتى
تموت أو تُقتل --- والعقبة: لا عامل يعيد تنشيطها
كلها، والخلايا المعاد تنشيطها لا تُقتل على نحو موثوق. والتحرير المورّثي:
استئصال \omterm{def:b3:virology:virus}{الفيروس} المندمج أو هدم المستقبل CCR5 في خلايا
المريض --- والعقبة: بلوغ كل خلية. (وقد شفى استبدالُ النقي بخلايا
متبرع معوزة CCR5 حفنةً من المرضى بخطر لا يُقبل
إلا حين تكون الزراعة لازمة أصلًا.)
\textbf{22.} $1 - 1/R_{0}$: أي \qty{67}{\%} من أجل $R_{0} = 3$، و\qty{80}{\%} من أجل
$5$.
\textbf{23.} العلاج يجعل $\beta \to 0$ داخل المريض، فيهبط الحمل
بحسب المبرهنة إلى لا شيء تقريبًا، ويتوقف الانتقال، وهو
متناسب مع الحمل؛ وعلاج كل مصاب يدفع $R_{0}$ الجمهري
دون الواحد.
\textbf{24.} $f = (1 - 1/R_{0})/E$: فعند $R_{0} = 4$ و$E = 0.7$، تكون $f =
0.75/0.7 = 1.07$ --- أي أكثر من الجميع: \omterm{def:b3:virology:drugs}{فاللقاح} وحده لا يستطيع بلوغ
العتبة؛ وعند $E = 0.9$، تكون $f = 0.83$.
\textbf{25.} نحو $10^{10}$ \omterm{def:b3:virology:virus}{فيريون} يُنتج في اليوم؛ وغير مكشوف في
اليوم $17$ تقريبًا؛ ونحو $4\times 10^{-6}$ طافر ثلاثي في السنة على العلاج.
\end{solution}
